% Generated by report/make_restated.py from theory/B1_preconditioned_saga.tex: do not edit (edit the note and rerun).
\begin{restated}{Remark}{rem:B1-memory-bottleneck}{B1}{0}{memory-bottleneck regime}
The one-step contraction factor is at most $1-\min\{\frac1{3n},\frac{3}{16\kappa'}\}$, with $\kappa'=\kappaP$
for the certified bound \eqref{eq:B1-onestep} and $\kappa'=R_P^2/\muEff^P(\theta_{t-1})$ for the state-dependent
bound \eqref{eq:B1-effective-rate} (a rate only under a deterministic bound on $\kappa'$ along the trajectory);
the memory term $\frac1{3n}$ is the binding one iff $\kappa'\le\frac{9n}{16}$. For
$P=\hat\Sigma+\lambda I$ and the trajectory constants of Corollary~\ref{cor:B1-kappa-traj}, a sufficient
condition is therefore
\begin{equation*}\tag{\ref*{eq:B1-bottleneck-condition}}
n\ \ge\ \frac{16}{9}\cdot\frac{d_\infty(\lambda)/4+1}{m}=\frac{4d_\infty(\lambda)+16}{9m}
\qquad\Big(\le\frac{4C\,d_{\mathrm{eff}}(\lambda)+16}{9m}\ \text{ under Assumption~\ref{ass:B1-coherence}}\Big),
\end{equation*}
with $m$ a deterministic constant such that $m(\theta_t)\ge m$ for all $t$ almost surely: a
trajectory-dependent $m$ cannot be pulled out of the expectations when the one-step bounds are iterated (see
the case of a non-constant metric in Section~\ref{sec:B1-cost}), and without an a priori confinement of the
iterates Proposition~\ref{prop:B3-region} of the B3 note, transferred, gives the bound only up to the exit time
from a region (end of Section~\ref{sec:B1-choice}). In this regime ($n\gg d_{\mathrm{eff}}/m$) the bound on
the contraction of $P$-SAGA is $1-\frac1{3n}$ whatever $\lambda$ and the spectrum of $\hat\Sigma$: it needs
$3\log(2n/\varepsilon)$ effective passes (the $\log(2n)$ coming from the prefactor $1+\frac n4$ of
\eqref{eq:B1-rate}), which is the $n\log(1/\varepsilon)$ term of the lower bound of
hypothesis (A1). That bound is a worst case over the problem class for span-based methods with i.i.d.\ sampling,
which cover $P$-SAGA only through Proposition~\ref{prop:B1-transfer} and for $P$ fixed independently of the
data; with full-gradient steps this term becomes of order $\frac{n}{1+\ln(n/\kappa)}\log\frac1\varepsilon$ when $\kappa\ll n$
(A1 note), and the accelerated complexity $O((n+\sqrt{n\kappa})\log(1/\varepsilon))$ of (A3) keeps it. For $P=I$ the
corresponding condition $n\ge\frac{16}{9}\frac{\max_i\norm{x_i}^2/4+\lambda}{m\sigma_d+\lambda}$ fails by a
factor of order $\max_i\norm{x_i}^2/(d_\infty(\lambda)\,\sigma_d)$ when $\lambda\ll m\sigma_d$; and with
global constants both conditions read $n\ge\frac{16}{9}\kappaP\ge\frac{16}{9}\kappa$ by
Proposition~\ref{prop:B1-nofreelunch}. This fails, by at least the factor $\frac{16}9$ and for every $P$, when
$\lambda=R^2/n$, i.e.\ $\kappa=n$, which is the setting of the experiment of Section~\ref{sec:B1-predictions}.
The experiments of \cite[Sec.~5.4.4]{bach2024learning} take $\mu=R_x^2/n$ with $R_x$ the bound on
$\norm{x_i}$: this is $\kappa=n$ in the book's normalization (smoothness constant $R_x^2$), but
$\kappa=\frac n4+1\le\frac{9n}{16}$ (for $n\ge4$) with the exact constant $R^2=R_x^2/4+\lambda$ used here, so
that there the condition holds for $P=I$ and the memory term is already the binding one with global constants
(in our reproduction of that setting, Figure~0 of the report with $d=40$, it also holds for
$P=\hat\Sigma+\lambda I$: $\kappaP\le347$ for $n=10^3$ and $\kappaP\le2814$ for $n=10^4$).
\end{restated}
